\documentclass[9pt,a4paper]{extarticle}
\usepackage{parskip}

% ======================
% Packages
% ======================
\usepackage[utf8]{inputenc}
\usepackage[T1]{fontenc}
\usepackage[brazil]{babel}
\usepackage{amsmath, amssymb}
\usepackage{amsthm}
\usepackage{geometry}
\usepackage{listings}
\usepackage{xcolor}
\usepackage{hyperref}
\usepackage{booktabs}
\usepackage{array}
\usepackage{tabularx}
\usepackage{enumitem}

\setlist{nosep}
\geometry{margin=2cm}

% ======================
% Listings setup (Python)
% ======================
\definecolor{codebg}{rgb}{0.95,0.95,0.95}
\lstset{
    language=Python,
    backgroundcolor=\color{codebg},
    basicstyle=\ttfamily\footnotesize,
    keywordstyle=\color{blue},
    commentstyle=\color{gray},
    stringstyle=\color{red!70!black},
    showstringspaces=false,
    frame=single,
    breaklines=true,
    inputencoding=utf8
}

% ======================
% Document
% ======================
\begin{document}

\title{Roteiro de Aula 2 -- Conceitos Fundamentais e Condições de Otimalidade}
\author{Disciplina: PCC179 -- Otimização Não Linear}
\date{\today}
\maketitle

\section{Mínimos e Máximos Globais}

\subsection*{Definição}
Seja $f:\mathbb{R}^n \to \mathbb{R}$ uma função.  
Um ponto $x^\star \in \mathbb{R}^n$ é dito:
\begin{itemize}
    \item um \textbf{mínimo global} se
    \[
    f(x^\star) \le f(x), \quad \forall x \in \mathbb{R}^n;
    \]
    \item um \textbf{máximo global} se
    \[
    f(x^\star) \ge f(x), \quad \forall x \in \mathbb{R}^n.
    \]
\end{itemize}

\paragraph{Interpretação:}
O ponto ótimo global é aquele que apresenta o melhor valor da função objetivo em todo o domínio — não apenas em torno de $x^\star$, mas em qualquer ponto possível.

\subsection*{Exemplo geométrico}
A função quadrática $f(x) = (x-2)^2 + 1$ tem um único mínimo global em $x^\star = 2$, pois $f(x) \ge 1$ para todo $x$.

\begin{lstlisting}
import numpy as np
import matplotlib.pyplot as plt

x = np.linspace(-1, 5, 400)
f = (x - 2)**2 + 1
plt.plot(x, f)
plt.scatter(2, 1, c='red', label='Minimo global (x=2)')
plt.title("f(x) = (x-2)^2 + 1")
plt.legend()
plt.xlabel("x"); plt.ylabel("f(x)")
plt.grid(True); plt.show()
\end{lstlisting}

\section{Mínimos e Máximos Locais}

\subsection*{Definição}
Um ponto $x^\star$ é um \textbf{mínimo local} de $f$ se existir $\epsilon > 0$ tal que:
\[
f(x^\star) \le f(x), \quad \forall x \in B_\epsilon(x^\star),
\]
onde $B_\epsilon(x^\star) = \{ x \in \mathbb{R}^n : \|x - x^\star\| < \epsilon \}$ é a vizinhança de raio $\epsilon$ em torno de $x^\star$.

Analogamente, $x^\star$ é um \textbf{máximo local} se:
\[
f(x^\star) \ge f(x), \quad \forall x \in B_\epsilon(x^\star).
\]

\subsection*{Observação}
- Um ponto pode ser mínimo local, mas não global.  
- Em funções não convexas, geralmente há múltiplos mínimos locais.  
- Se $f$ é convexa, todo mínimo local é também global.

\paragraph{Exemplo}
Considere:
\[
f(x) = x^4 - 3x^2 + 2.
\]
Esta função tem dois mínimos locais e um máximo local.

\begin{lstlisting}
x = np.linspace(-3, 3, 400)
f = x**4 - 3*x**2 + 2
plt.plot(x, f)
plt.title("f(x) = x^4 - 3x^2 + 2 (multiplos minimos locais)")
plt.xlabel("x"); plt.ylabel("f(x)")
plt.grid(True)
plt.show()
\end{lstlisting}

\section{Condições de Otimalidade de Primeira Ordem}

\subsection*{Teorema (Condição Necessária de Primeira Ordem)}

Se $f:\mathbb{R}^n \to \mathbb{R}$ é diferenciável e $x^\star$ é um ponto de mínimo ou máximo local, então:
\[
\nabla f(x^\star) = 0.
\]
Esse ponto $x^\star$ é chamado de \textbf{ponto crítico} ou \textbf{estacionário}.

\subsection*{Interpretação Geométrica}
No ponto crítico, a taxa de variação direcional de $f$ em qualquer direção $d$ é zero:
\[
\nabla f(x^\star)^\top d = 0, \quad \forall d.
\]
Ou seja, a tangente do gráfico é horizontal (em 1D) ou o gradiente é o vetor nulo (em $n$D).

\subsection*{Exemplo}
Para $f(x) = x^3 - 3x$, temos:
\[
f'(x) = 3x^2 - 3 = 3(x^2 - 1) \implies f'(x)=0 \Rightarrow x = \pm 1.
\]

\begin{lstlisting}
x = np.linspace(-2, 2, 400)
f = x**3 - 3*x
plt.plot(x, f)
plt.axvline(0, color='k', linewidth=0.8)
plt.scatter([-1,1], [-2,2], c='red', label='Pontos criticos')
plt.legend(); plt.grid(True)
plt.title("Condicao de 1a ordem: f'(x)=0 em x = +-1")
plt.show()
\end{lstlisting}

\section{Condições de Otimalidade de Segunda Ordem}

\subsection*{Teorema (Condição Suficiente de Segunda Ordem)}
Suponha que $f:\mathbb{R}^n \to \mathbb{R}$ seja duas vezes diferenciável e que $\nabla f(x^\star) = 0$. Então:
\begin{itemize}
    \item Se $\nabla^2 f(x^\star)$ é \textbf{definida positiva}, $x^\star$ é um \textbf{mínimo local estrito}.
    \item Se $\nabla^2 f(x^\star)$ é \textbf{definida negativa}, $x^\star$ é um \textbf{máximo local estrito}.
    \item Se $\nabla^2 f(x^\star)$ é \textbf{indefinida}, $x^\star$ é um \textbf{ponto de sela}.
\end{itemize}

\paragraph{Em 1 dimensão:}
\[
f''(x^\star) > 0 \Rightarrow \text{mínimo local}, \quad
f''(x^\star) < 0 \Rightarrow \text{máximo local}.
\]

\subsection*{Exemplo}
Para $f(x) = x^3 - 3x$, temos:
\[
f''(x) = 6x.
\]
Logo:
\[
f''(-1) = -6 < 0 \Rightarrow \text{máximo local}, \quad
f''(1) = 6 > 0 \Rightarrow \text{mínimo local}.
\]

\subsection*{Visualização 2D (ponto de sela)}
A função $f(x,y) = x^2 - y^2$ tem gradiente nulo em $(0,0)$, mas Hessiana indefinida:
\[
\nabla^2 f(0,0) =
\begin{bmatrix}
2 & 0 \\ 0 & -2
\end{bmatrix}.
\]
Logo, $(0,0)$ é um ponto de sela.

\begin{lstlisting}
from mpl_toolkits.mplot3d import Axes3D
X, Y = np.meshgrid(np.linspace(-2,2,100), np.linspace(-2,2,100))
Z = X**2 - Y**2

fig = plt.figure()
ax = fig.add_subplot(111, projection='3d')
ax.plot_surface(X, Y, Z, cmap='coolwarm', alpha=0.8)
ax.scatter(0, 0, 0, c='k', s=40)
ax.set_title("f(x,y) = x^2 - y^2 (ponto de sela)")
plt.show()
\end{lstlisting}

\pagebreak

\section{Resumo das Condições de Otimalidade}

\begin{table}[!ht]
\centering
\small
\setlength{\tabcolsep}{5pt}
\renewcommand{\arraystretch}{1.15}
\begin{tabular}{@{}lcl@{}}
\toprule
\textbf{Condição} & \textbf{Formulação} & \textbf{Interpretação} \\
\midrule
1ª Ordem (necessária) & $\nabla f(x^\star) = 0$ & Nenhuma direção reduz $f$ localmente. \\
2ª Ordem (suficiente) & $\nabla^2 f(x^\star)$ definida positiva & Mínimo local estrito. \\
2ª Ordem (suficiente) & $\nabla^2 f(x^\star)$ definida negativa & Máximo local estrito. \\
Ponto de sela & $\nabla^2 f(x^\star)$ indefinida & Nem mínimo nem máximo. \\
\bottomrule
\end{tabular}
\caption{Resumo das condições de otimalidade (problemas sem restrições).}
\end{table}


\section{Expansão em Série de Taylor}

A expansão em série de Taylor é uma ferramenta fundamental em otimização, pois permite aproximar funções diferenciáveis por polinômios, revelando o comportamento local de $f$ em torno de um ponto $x^\star$. Essa aproximação é a base teórica das condições de otimalidade e de muitos algoritmos (como métodos de Newton).

\subsection*{Definição (Série de Taylor de Primeira e Segunda Ordem)}

Seja $f:\mathbb{R}^n \to \mathbb{R}$ uma função com derivadas contínuas até segunda ordem.  
A expansão de $f$ em torno de um ponto $x^\star$ é dada por:

\[
f(x) = f(x^\star) 
+ \nabla f(x^\star)^\top (x - x^\star)
+ \tfrac{1}{2}(x - x^\star)^\top \nabla^2 f(x^\star) (x - x^\star)
+ o(\|x - x^\star\|^2),
\]

onde:
\begin{itemize}
    \item $\nabla f(x^\star)$ é o \textbf{vetor gradiente} em $x^\star$;
    \item $\nabla^2 f(x^\star)$ é a \textbf{matriz Hessiana} em $x^\star$;
    \item $o(\|x - x^\star\|^2)$ representa termos de ordem superior que se tornam desprezíveis para $x$ próximo de $x^\star$.
\end{itemize}

\subsection*{Interpretação Geométrica}

- O termo \textbf{linear} $\nabla f(x^\star)^\top (x - x^\star)$ representa o plano tangente de $f$ em $x^\star$.  
- O termo \textbf{quadrático} $\tfrac{1}{2}(x - x^\star)^\top \nabla^2 f(x^\star)(x - x^\star)$ descreve a curvatura local.  
- Em 1D, temos:
  \[
  f(x) \approx f(x^\star) + f'(x^\star)(x - x^\star) + \tfrac{1}{2}f''(x^\star)(x - x^\star)^2.
  \]

\paragraph{Observação:}
A partir dessa expansão, podemos entender:
\begin{itemize}
    \item Se o termo linear é nulo ($\nabla f(x^\star)=0$), o ponto é estacionário.
    \item O termo quadrático (Hessiana) determina se esse ponto é mínimo, máximo ou sela.
\end{itemize}

\subsection*{Exemplo Numérico (1D)}
Considere $f(x) = e^x$ e $x^\star = 0$.  
A expansão de Taylor até segunda ordem é:
\[
f(x) \approx f(0) + f'(0)x + \tfrac{1}{2} f''(0)x^2 = 1 + x + \tfrac{x^2}{2}.
\]

\begin{lstlisting}
import numpy as np
import matplotlib.pyplot as plt

x = np.linspace(-2, 2, 400)
f = np.exp(x)
taylor = 1 + x + 0.5*x**2

plt.plot(x, f, label='f(x) = e^x')
plt.plot(x, taylor, '--', label='Aproximacao de Taylor (2a ordem)')
plt.legend(); plt.grid(True)
plt.title("Expansao de Taylor de f(x)=e^x em torno de x=0")
plt.xlabel("x"); plt.ylabel("f(x)")
plt.show()
\end{lstlisting}

\subsection*{Exemplo Multidimensional}
Para $f(x,y) = x^2 + xy + y^2$, em torno de $(0,0)$:
\[
\nabla f(0,0) =
\begin{bmatrix} 0 \\ 0 \end{bmatrix},
\quad
\nabla^2 f(0,0) =
\begin{bmatrix}
2 & 1 \\
1 & 2
\end{bmatrix}.
\]
Logo, a expansão até segunda ordem é:
\[
f(x,y) \approx x^2 + xy + y^2.
\]
Nesse caso, a função é puramente quadrática e a expansão coincide exatamente com a função original.

\subsection*{Resumo Intuitivo}
\begin{itemize}
    \item A série de Taylor permite \textbf{linearizar} e \textbf{aproximar} $f$ localmente.
    \item Em algoritmos de otimização, essa aproximação serve para definir a direção e o tamanho do passo.
    \item As condições de primeira e segunda ordem derivam diretamente dessa expansão.
\end{itemize}

\section*{Referências}
\begin{itemize}
    \item Beck, A. (2017). \textit{Introduction to Nonlinear Optimization: Theory, Algorithms, and Applications with MATLAB}, 2nd Edition. Springer.
    \item Nocedal, J., Wright, S. J. (2006). \textit{Numerical Optimization}. Springer.
\end{itemize}


\end{document}
